\section{Closing the common-rate 64-cycle game}
\label{sec:cycle64-endgame}

A global upper bound can be obtained from one adverse labeling if every support and every rate vector on that labeling is covered. The census on word \(2549776587809433532\) shows that only two admissible supports have a closed-box maximum above \(0.10118\). This threshold falls \emph{below} both constructive guarantees. Consequently the only unresolved task on that word is an upper bound for arbitrary rates on two particular finite supports. Each support has a single preferred leaf in its sparse half. The resulting mass constraint is sufficiently strong to close the gap.

Throughout this section a rule can read the labeling before returning its support and rates. Its response to a given labeling must be usable for every \(0<\varepsilon<1/52\). Under this convention the optimum for the literal prefactor certificate lies in the interval
\[
 10.1190251762367
 <100\,V_{\rm literal}<
 10.1190251762403 .
\]
More precisely we have an exact implicit value and a real-valued attaining policy. The interval measures the numerical evaluation of that value. The original 10.0046 percent rule used the same winning support decisions with other rates. Once these rates are optimized, the ability to change the occupied residues gives no further \emph{worst-label} gain under the stated certificate. An individual label can have better layouts. The matching upper uses the rule's obligation to handle the particular labeling in the census.

\subsection{Why exactly two calculations remain}

Record a mask by the integer with bit \(t\) set when residue \(t\) is nominally occupied. Denote the surviving masks by
\[
 S_1=1230983613671479063,\qquad S_2=1518083709116552983 .
\]
Either can retain points in 28 residue classes. On both, the odd half of the root contains \(1,9,17,21,27,33\). Only \(27\) lies in its preferred child. Write \(b=a_{27}\). The residue \(46\) in the even half will absorb the cardinality constraint. It has coefficients \((p_{46},q_{46})=(0,33)\), and the odd anchor has \((p_{27},q_{27})=(0,65)\). Here \(P=\sum_t p_t a_t,\ Q=\sum_t q_t a_t\) include the finer-level contribution through the unit term in each \(q_t\).

Let \(I=S_i\setminus\{27,46\}\) and \(J=\{1,9,17,21,33\}\). To reproduce the coefficients, at each node form all directed pair differences between the supported residues, and count the missing odd increments. Each odd child leaf acquires twice that count in \(p_t\) and the node width in \(q_t\). Simple summation gives
\[
 \begin{array}{c|c|c|c}
 & |I| & \sum_{t\in I}p_t & D_i=\sum_{t\in I}(q_t-33)\\ \hline
 S_1&26&116&262\\
 S_2&26&116&254
 \end{array}
\]
The difference of eight reflects the dyadic exposures of three even-half leaves exchanged between the two supports. These sums alone do not establish which support is worse. An allocation could give distinct weights to all 26 sites. The following coefficient inequality controls that possibility.

\subsection{The disappearing mass of the five sparse leaves}

Let \(C\) be a common cap on all leaf rates and put \(Y=\sum_{t\in J}(C-a_t)\). The sparse half has total mass \(b+5C-Y\) and preferred mass \(b\). For a rate vector valid at every exponent, its limiting mark requires
\[
 b\ \ge\ 0.32\exp\{(b+5C-Y)/16\}.
 \tag{E1}
\]
For \(s\in[0,5C]\) let \(B(s)\) be the small positive solution of \(B=0.32\exp\{(B+s)/16\}\), and write \(B=B(5C)\). The function \(b\mapsto b\exp(-b/16)\) is strictly increasing up to \(16\). All relevant \(b\) are below \(C<2.4\), so (E1) implies \(b\ge B(5C-Y)\). Implicit differentiation gives \(B'(s)=B(s)/(16-B(s))\). At both caps used below, the exact enclosures give \(B(5C)<0.71\), hence \(0<B'(s)<1/15\) throughout the interval. Integration yields
\[
 B-b\ \le\ Y/15.
 \tag{E2}
\]
This is the quantitative reason a rate reduction on the five sparse leaves cannot pay for moving appreciable mass away from their anchor. It also covers the case where the anchor already receives more than \(B\).

For \(i=1,2\) define
\[
 z_i(C)=\frac{116C}{2112+D_i C+32B},\qquad
 v_t=p_t-z_i(C)(q_t-33).
\]
The hypothetical allocation that fills all 26 sites in \(I\) at \(C\), puts \(B\) on \(27\), and sends the remainder to \(46\) has \(P/Q=z_i(C)\). To compare an arbitrary allocation to it, eliminate \(a_{46}=64-b-\sum_{I}a_t\). Since \(64\cdot33=2112\), subtraction gives the exact identity
\[
 P-z_i(C)Q
   =32z_i(C)(B-b)-\sum_{t\in I}v_t(C-a_t).
 \tag{E3}
\]
The bound \(0.10\le z_i(C)\le0.11\) contains both critical ratios for both versions of the mark. For the two physical supports, direct integer coefficients give \(v_t>0\) throughout this range and \(v_t>32z_i(C)/15\) for \(t\in J\). Each difference is affine in \(z\); it suffices to check the two rational endpoints. Paying the first term of (E3) from the five sparse deficits using (E2) gives \(P-z_i(C)Q\le0\). Since \(Q>0\), the candidate is an upper bound valid for every rate allocation on this support. This is the weighted finite inequality covered by the general kernel lemma in this study, after supplying its computed coefficients and the scalar analytic bound (E2).

We can now compare supports without any assumption of equal rates in the competing profiles. The first support has a larger \(D_i\) and identical \(C,B\), hence \(z_1(C)<z_2(C)\). All other supports have maximum at most \(0.10118<z_2(C)\) by the census. Masks with no allowable preferred marker or too few positive root preferred leaves have already been excluded by necessary conditions. In particular, a search that listed only 27-leaf extremal vertices would not have justified discarding the 28-leaf nominal supports in this calculation. A positive anchor that disappears at a capped-simplex vertex can determine the true optimum over its parent support.

\subsection{The two sharp constants and the distinction between maximum and supremum}

Under the literal-prefactor criterion the homogeneous-leaf cap \(C_*\) is the larger root of \(10 e^{-C_*}\sqrt{C_*/2}=1\). Under the standard strict product criterion its cap \(C_0\) is the larger root of \(10 e^{-C_0}(C_0/2)^{5190/10381}=1\). The latter exponent is the smallest value of \(3r/(6r+1)\) in the admissible source range and is realized there. In either case write \(B_{\circ}(C)=B(5C)\) in the notation above. The census and projection give
\[
 V_{\rm literal}=\frac{116C_*}{2112+254C_*+32B_{\circ}(C_*)},\qquad
 V_{\rm standard}=\frac{116C_0}{2112+254C_0+32B_{\circ}(C_0)}.
\]
The second value is interpreted as a supremum. Its verified percentage interval is
\(10.118989577495<100V_{\rm standard}<10.118989577849\).

For the lower directions, use the policy that kept the original support decisions. The ordinary decisions exceed the higher of the two critical fractions at interior rational rates. In the exceptional decision take 26 common productive rates \(C\), one anchor rate \(B_{\circ}(C)\), and one residual rate \(64-26C-B_{\circ}(C)\). Of its 23 distinct cell types, 21 have a uniform positive logarithmic surplus on a rational rectangle containing the relevant corner; the smallest certified surplus exceeds \(0.1473\) in both evaluations. The two remaining types are the homogeneous high-rate leaf and the marked sparse half. In the prefactor case each has \(10e^{-\mu}\sqrt\nu=1\) at its exact corner. The powered left side of the literal comparison is therefore exactly \(1/\sqrt\nu\). For both relevant values of \(\nu\), verified rational bounds give \(\nu(54/67)^2<1\), while \((h+2)/(h+3)<54/67\) for every admissible \(\varepsilon\). Both literal inequalities remain strict. Thus the real rates at \(C_*,B_{\circ}(C_*)\) actually attain the common-rate guarantee. Fixed-count samples on increasingly late arithmetic hosts implement these rates asymptotically.

In the standard case the homogeneous leaf at \(C_0\) gives equality at the smallest allowed tile exponent. Taking \(C<C_0\) tending upward to \(C_0\), and choosing the anchor just above its small-root boundary, gives strict inequalities and limits to the claimed ratio. Conversely, equality in (E3) requires all 26 productive sites to be filled to their cap, because the coefficient inequalities have strict surplus. Those homogeneous sites cannot take the exact cap under the standard strict test. This explains why two numerically adjacent constants have different attainment statements.

The logical conclusion concerns a finite decision problem and its specified passage to arithmetic families. It uses the existence of compatible selected masks. Nothing in the matching lower and upper values identifies the set returned by a particular classical selection in the Lean source. A different proof of the original score, a finer quotient, or profiles whose rates vary with the source exponent define different optimization problems.
